% Exportado desde agregados; no editar cifras manualmente.
\begin{table}[!htbp]
\centering
\caption{Recepción propia y estados de información por ronda, directas 60+, 2001--2021}\label{tab:tesis-cobertura}
\footnotesize\setlength{\tabcolsep}{3pt}\renewcommand{\arraystretch}{1.25}
\begin{tabularx}{\linewidth}{p{.075\linewidth}p{.085\linewidth}YYYp{.135\linewidth}p{.13\linewidth}}
\toprule
\textbf{Ronda} & \textbf{N elegible} & \textbf{K: \% (N)} & \textbf{A: \% (N)} & \textbf{Unión: \% (N)} & \textbf{ND K/A/unión} & \textbf{Ambas negativas} \\
\midrule
2001 & 6\,812 & \ensuremath{24,74}\ \allowbreak\ensuremath{(6\,790)} & \ensuremath{9,81}\ \allowbreak\ensuremath{(6\,789)} & \ensuremath{33,23}\ \allowbreak\ensuremath{(6\,789)} & 22/\allowbreak{}23/\allowbreak{}23 & 4\,533 \\
2003 & 6\,970 & \ensuremath{25,53}\ \allowbreak\ensuremath{(6\,894)} & \ensuremath{12,98}\ \allowbreak\ensuremath{(6\,894)} & \ensuremath{36,64}\ \allowbreak\ensuremath{(6\,894)} & 76/\allowbreak{}76/\allowbreak{}76 & 4\,368 \\
2012 & 8\,888 & \ensuremath{32,26}\ \allowbreak\ensuremath{(8\,738)} & \ensuremath{17,99}\ \allowbreak\ensuremath{(8\,750)} & \ensuremath{47,42}\ \allowbreak\ensuremath{(8\,734)} & 150/\allowbreak{}138/\allowbreak{}154 & 4\,592 \\
2015 & 9\,640 & \ensuremath{43,76}\ \allowbreak\ensuremath{(9\,504)} & \ensuremath{27,65}\ \allowbreak\ensuremath{(9\,513)} & \ensuremath{65,79}\ \allowbreak\ensuremath{(9\,508)} & 136/\allowbreak{}127/\allowbreak{}132 & 3\,253 \\
2018 & 9\,107 & \ensuremath{35,89}\ \allowbreak\ensuremath{(9\,047)} & \ensuremath{30,56}\ \allowbreak\ensuremath{(9\,045)} & \ensuremath{62,74}\ \allowbreak\ensuremath{(9\,039)} & 60/\allowbreak{}62/\allowbreak{}68 & 3\,368 \\
2021 & 9\,424 & \ensuremath{34,82}\ \allowbreak\ensuremath{(9\,366)} & \ensuremath{52,59}\ \allowbreak\ensuremath{(9\,363)} & \ensuremath{67,97}\ \allowbreak\ensuremath{(9\,366)} & 58/\allowbreak{}61/\allowbreak{}58 & 3\,000 \\
\bottomrule
\end{tabularx}
\notacuadro{Cada tasa usa el N conocido entre paréntesis; ND se cuenta sobre el N elegible de la fila y no significa negativa. Unión: alguna afirmativa; no: ambas negativas. K y A se superponen. Las filas son cortes separados, no una serie homogénea a través de rupturas K 2003--2012/K 2015--2018 o A 2001--2003/A 2018--2021.}
\fuentecuadro{Elaboración propia, ENASEM; agregados de notebooks 06--08 y validación focalizada del informe general, secciones 10.2--10.7. Sin ponderación nacional.}
\end{table}
